\inicializarSubseccion{Analisis de la primera pregunta (Por parte de Viviana)}
\textbf{1. ¿Cuáles de estas combinaciones te gustarían en el producto? Selecciona por lo menos dos opciones.}

Para esta pregunta se formaron 4 conjuntos ya que había 4 diferentes opciones los cuales los definimos como $A,B,C,D$:

\begin{itemize}
    \item $A$ = adaptador + preservar el calor o frío con 87
    \item $B$ = adaptador + compartimento secreto con 39
    \item $C$ = adaptador + material ecológico con 40
    \item $D$ = adaptador + antideslizante con 41
\end{itemize}

$A,B,C,D$, serían contenidos en un conjunto universal $U$ que sería el total de personas encuestadas, el cual es de 100 personas.:

\begin{itemize}
    \item $A = 87$
    \item $B = 39$
    \item $C = 40$
    \item $D = 41$
\end{itemize}

Dándonos que en el conjunto de $A\cap B$, $A\cap C$, $A\cap D$, $B\cap C$, $B\cap D$, $C\cap D$ nunca vayan a estar vacios, ya que es necesario que se escojan por lo menos dos opciones y, a su vez, podrían existir triples intercepciónes entre $A\cap B\cap C$, $A\cap B\cap D$, $A\cap C\cap D$, $B\cap C\cap D$ y una intercepción entre los cuatro conjuntos $A\cap B\cap C\cap D$.

Realizando la cuenta en porcentaje, podemos saber que $P(A) = \frac{87}{100} = 0.87$ o $87\%$.

Gracias a esto podemos ver que la gran mayoria de nuestro mercado quiere el adaptador que pueda preservar el calor o el frio.
